\documentclass{article}
\usepackage{amsmath, amssymb, amsthm}

\setlength{\parindent}{0pt}

\newtheorem{claim}{Claim}
\newtheorem*{corollary}{Corollary}

\begin{document}

\begin{claim}
    Let $n = p \cdot q$ where $p, q$ are both odd primes and $x = e \cdot d - 1$ where $e$ is the public key and
    $d$ is the private key of the RSA protocol. Then $\phi(n) \mid x$.
\end{claim}

\begin{proof}
    By definition of the RSA private key, $d$ is the modular inverse of $e$ modulo $\phi(n)$. Hence:
    \[
    ed \equiv 1 \pmod{\phi(n)}
    \]
    Then it follows that:
    \[
    \phi(n) \mid (ed - 1)
    \]
    Since $x = e \cdot d - 1$, substituting in gives us:
    \[
    \phi(n) \mid x
    \]

    \medskip

    Hence $\phi(n) \mid x$.
\end{proof}

\begin{claim}
    $4 \mid x$
\end{claim}

\begin{proof}
    By definition of $\phi(n)$:
    \[
    \phi(n) = (p - 1)(q - 1)
    \]
    Since $p$ and $q$ are odd, $p - 1$ and $q - 1$ are even. Therefore $2 \mid p - 1$ and $2 \mid q - 1$. Hence:
    \[
    4 = 2 \cdot 2 \mid (p - 1)(q - 1) = \phi(n)
    \]
    Since $4 \mid \phi(n)$ and, by Claim 1, $\phi(n) \mid x$, it follows by transitivity of divisibility 
    that $4 \mid x$.

    \medskip

    Hence $4 \mid x$.
\end{proof}

\begin{claim}
    If $\gcd(r, n) = 1$, then $r^{x} \equiv 1 \pmod{n}$.
\end{claim}

\begin{proof}
    By Claim 1, $\phi(n) \mid x$. Hence by the definition of divisibility there exists an integer $k$ such that:
    \[
    x = k \cdot \phi(n)
    \]
    Since $\gcd(r, n) = 1$, by Euler's theorem:
    \[
    r^{\phi(n)} \equiv 1 \pmod{n}
    \]
    Raising both sides to the power $k$:
    \[
    \begin{aligned}
    (r^{\phi(n)})^{k} &\equiv 1^{k} \pmod{n} \\
    r^{k \cdot \phi(n)} &\equiv 1 \pmod{n} \\
    \end{aligned}
    \]
    Since $x = k \cdot \phi(n)$,
    \[
    r^{x} \equiv 1 \pmod n.
    \]

    \medskip

    Hence if $\gcd(r, n) = 1$, then $r^{x} \equiv 1 \pmod{n}$.
\end{proof}

\begin{claim}
    If $y$ is a nontrivial square root of $1 \pmod{n}$, then you can factor $n$.
\end{claim}

\begin{proof}
    By definition of $y$:
    \[
    y^{2} \equiv 1 \pmod{n}
    \]
    Thus:
    \[
    n \mid (y^{2} - 1)
    \]
    Since $n = pq$ from Claim 1, substituting in gives:
    \[
    pq \mid (y^{2} - 1)
    \]
    Factoring $y^{2} - 1$ gives:
    \[
    pq \mid (y - 1)(y + 1)
    \]

    \medskip

    Since $pq \mid (y - 1)(y + 1)$, $p \mid (y - 1)(y + 1)$ and $q \mid (y - 1)(y + 1)$. Since $p$ and $q$ are prime, 
    by Euclid's lemma:
    \[
    p \mid (y - 1) \quad \text{or} \quad p \mid (y + 1)
    \]
    Similarly:
    \[
    q \mid (y - 1) \quad \text{or} \quad q \mid (y + 1)
    \]
    
    \medskip

    Suppose $p \mid (y - 1)$ and $q \mid (y - 1)$. Since $p$ and $q$ are distinct primes:
    \[
    pq \mid (y - 1)
    \]
    As $n = pq$, it follows that:
    \[
    n \mid (y - 1)
    \]
    Hence:
    \[
    y \equiv 1 \pmod{n}
    \]
    However, this contradicts the assumption that $y$ is a nontrivial square root of $1 \pmod{n}$. Therefore this 
    case is impossible.

    \medskip

    Similarly, if $p \mid (y + 1)$ and $q \mid (y + 1)$, then:
    \[
    pq \mid (y + 1)
    \]
    Then:
    \[
    n \mid (y + 1)
    \]
    Thus:
    \[
    y \equiv -1 \pmod{n}
    \]
    Which again, contradicts the assumption that $y$ is a nontrivial square root of $1 \pmod{n}$. Therefore this 
    case is also impossible.

    \medskip

    Therefore, exactly one of $p$ and $q$ divides $y - 1$ and the other divides $y + 1$. Without loss of generality, 
    suppose $p \mid (y - 1)$ and $q \mid (y + 1)$. Since $p \mid (y - 1)$, $p \mid n$, and $q \nmid (y - 1)$, 
    the only prime divisor of $n$ that divides $y - 1$ is $p$. Hence:
    \[
    \gcd(y - 1, n) = p
    \]
    Likewise for $q$:
    \[
    \gcd(y + 1, n) = q
    \]

    \medskip

    Therefore if $y$ is a nontrivial square root of $1 \pmod{n}$, you can factor $n$.
\end{proof}

\begin{corollary}
    If, in addition to knowing $n$ and the public key $e$, you also knew the private key $d$, then you can factor $n$.
\end{corollary}

\begin{proof}
    By the given statement, at least half of the admissable choices of $r$ produce a nontrivial square root of $1 
    \pmod{n}$. By Claim 4, any such square root yields a factorisation of $n$. Hence repeated random choices of $r$ 
    eventually factor $n$.
\end{proof}

\end{document}
